\section{Week 3: Tokenization and Embeddings — From Text to Vectors}

\subsection{Learning Objectives}

By the end of this week, you should be able to:
\begin{itemize}
  \item Explain why tokenization is a design choice, not a neutral preprocessing step
  \item Compare word-level and subword-level tokenization
  \item Understand embeddings as learned geometric representations
  \item Reason about how meaning, similarity, and analogy emerge in vector spaces
  \item Identify system-level consequences of tokenization decisions
\end{itemize}

\subsection{LLM Components Covered}

\textbf{Input representation}
\begin{itemize}
  \item Tokenization (word vs subword)
  \item Vocabulary construction
  \item Token IDs and special tokens
\end{itemize}

\textbf{Representation learning}
\begin{itemize}
  \item Embedding layers
  \item Semantic geometry
  \item Similarity and clustering
\end{itemize}

\subsection{Conceptual Core: Language as Discrete Symbols}

Neural networks operate on vectors in $\mathbb{R}^n$, but natural language is discrete and symbolic.
Tokenization is the process that defines the mapping:

\[
\text{text} \rightarrow \text{tokens} \rightarrow \text{integers} \rightarrow \text{vectors}
\]

This mapping is:
\begin{itemize}
  \item lossy,
  \item irreversible,
  \item and deeply consequential.
\end{itemize}

Once tokenized, the model never ``sees'' characters or words again — only token IDs.

\subsection{Word-Level Tokenization}

In word-level tokenization, text is split on whitespace or punctuation.
This approach is intuitive but brittle:
\begin{itemize}
  \item out-of-vocabulary (OOV) words,
  \item poor handling of morphology,
  \item exploding vocabulary size.
\end{itemize}

Early neural language models relied heavily on this approach.

\subsection{Subword Tokenization}

Modern LLMs use subword tokenization (e.g.\ Byte Pair Encoding).

Key idea:
\begin{itemize}
  \item frequent character sequences become tokens,
  \item rare words are decomposed into smaller units.
\end{itemize}

This allows:
\begin{itemize}
  \item open-vocabulary modeling,
  \item graceful handling of neologisms,
  \item shared statistical strength across words.
\end{itemize}

\subsection{Embeddings as Learned Geometry}

Each token ID is mapped to a vector via an embedding matrix:

\[
\mathbf{e}_i \in \mathbb{R}^d
\]

These vectors are learned jointly with the rest of the model.
Distances and directions in embedding space capture statistical regularities of language.

Similarity is often measured via cosine similarity:

\[
\cos(\theta) = \frac{\mathbf{e}_i \cdot \mathbf{e}_j}{\|\mathbf{e}_i\| \|\mathbf{e}_j\|}
\]

\subsection{What Embeddings Encode (and What They Do Not)}

Embeddings encode:
\begin{itemize}
  \item distributional similarity,
  \item syntactic and semantic regularities,
  \item task-relevant biases from training data.
\end{itemize}

They do not encode:
\begin{itemize}
  \item ground truth,
  \item symbolic meaning,
  \item explicit ontologies.
\end{itemize}

\subsection{Systems Engineering Implications}

\textbf{Requirements}
\begin{itemize}
  \item Tokenization must be stable across training and inference
  \item Vocabulary changes imply retraining or incompatibility
\end{itemize}

\textbf{Architecture}
\begin{itemize}
  \item Tokenizers are part of the model contract
  \item Embedding tables dominate memory footprint at scale
\end{itemize}

\textbf{Verification \& Validation}
\begin{itemize}
  \item Check tokenization invariants
  \item Inspect embedding drift across versions
\end{itemize}

\textbf{Operations}
\begin{itemize}
  \item Token-level monitoring
  \item Input distribution shift detection
\end{itemize}

\subsection{Human Factors and Failure Modes}

Common failures include:
\begin{itemize}
  \item assuming ``word meaning'' is preserved,
  \item misunderstanding why similar prompts tokenize differently,
  \item prompt brittleness due to token boundaries.
\end{itemize}

These effects are invisible without explicit inspection.

\subsection{Key References}

\begin{itemize}
  \item \citet{mikolov2013efficient} — word embeddings
  \item \citet{sennrich2016bpe} — subword tokenization for neural MT
\end{itemize}

\subsection{Recommended University Lectures}

\begin{itemize}
  \item Stanford CS224N — word vectors and embeddings  
  \url{https://web.stanford.edu/class/cs224n/}

  \item MIT 6.864 — statistical NLP foundations  
  \url{https://nlp.csail.mit.edu}
\end{itemize}

\subsection{Practitioner Lab}

See \texttt{labs/week03/week03\_lab\_tokenization\_embeddings.ipynb}.

\subsection{Week 3 Takeaway}

\begin{quote}
Tokenization defines the language a model can speak; embeddings define the geometry in which it thinks.
\end{quote}

